\textbf{Task completion hides much of the damage caused by executed cover.}
With Qwen3:8B, live-state cover completed 102/120 requested tasks, but only
7/120 final states matched the intended business result. With Llama3.1:8B,
the corresponding counts were 76/120 and 3/120. Thus 95 and 73 completed
tasks still had incorrect global state. Relative to no cover, exact state
integrity fell from 116 to 7 and from 69 to 3; both paired reductions
survived Holm correction. Goal completion fell from 114 to 102 and from
82 to 76, with only the Qwen reduction significant after correction.
These constructed adverse writes make the practical consequence measurable:
an executed-cover evaluation needs a state-integrity endpoint alongside
task completion and traffic inference.

\textbf{The state loss bought no distinctive improvement in the measured
recognition scores.}
With one or two randomly placed decoys, the validation-selected LR/RF
attacker scored 61--63/120 across Qwen cover conditions and 72--82/120
across Llama cover conditions, compared with 118/120 and 113/120 without
cover. All ten cover-versus-baseline primary comparisons survived the
predeclared Holm correction. Live writes scored 62/120 and 80/120, within
the range already observed for proxy-discard, read, and isolated cover.
Residual exposure remained substantial: a last-destination rule reached
75--77/120 in Qwen cover conditions and exceeded the selected LR/RF attack.
These are observed attack results, not equivalence or privacy upper bounds.
Call-byte costs were 1.28--1.35 times the paired no-cover runs.

\textbf{Aggregate scores hide variation in the visible-response control's
individual execution paths.}
The hidden and visible isolated conditions returned identical cover
responses in all 120 pairs per model. Their honest action sequences
differed in 16 Qwen pairs and 65 Llama pairs. Mean prompt tokens rose by
18.1\% and 18.9\%, respectively. The predeclared visible-versus-hidden
comparisons for recognition, goal completion, and state integrity all had
Holm-adjusted $p=1$. The paired path differences therefore describe planner
behavior; they do not establish that exposing marked responses improves
or degrades average utility. Hidden-arm path variation is also retained
in the released audit.
